\cotefile{156-8}
\scanpage{50}
\mspage{41}
% archive sheet number: 50

$(\mathrm{i})\Longrightarrow(\mathrm{ii})$ Je \unsure{note} : prouvons la formule
\[
X_i\cap X_j=\bigcup_{k\leq i,j}X_k
\]
(côté $\supset$ est évident). Mais par $3'$ on a
\[
X_i=\bigcup_{i'\leq i}X_{i'}^\circ,\quad X_j=\bigcup_{j'\leq j}X_{j'}^\circ
\]
d'où
\[
X_i\cap X_j=\bigcup_{\substack{i'\leq i\\j'\leq j}}X_{i'}^\circ\cap X_{j'}^\circ
\]
Or par $2')$ on a $X_{i'}^\circ\cap X_{j'}^\circ=\varnothing$ sauf si $i'\ne j'$, d'où
\[
X_i\cap X_j=\bigcup_{k\leq i,j}X_k^\circ\subset\bigcup_{k\leq i,j}X_k\qquad\text{qfd}
\]

$(\mathrm{ii})\Longrightarrow(\mathrm{i})$. Cela résulte du

\deleted{\uline{Corollaire}} \inserted{Lemme} Supposons qu'on ait, $\forall i,j\in I$
\[
X_i\cap X_j\unsure{\supset|=}\bigcup_{k\leq i,j}X_k.
\]
Alors on a les propriétés :
\[
\forall i,j\in I,\ i\ne j\quad\deleted{\unsure{X_j}}\quad X_i^\circ\cap X_j^\circ=0\tag*{$2')$}
\]
\inserted{Si de plus $\bigcup_{i\in I}X_i^\circ=|\unsure{\Phi|P}|$ on a aussi}
\[
\forall i\in I\qquad X_i=\bigcup_{j\leq i}X_j^\circ.\tag*{\unsure{$3')$}}
\]

Prouvons $2')$ on a \inserted{\unsure{car|or}}
\[
X_i^\circ\cap X_j^\circ\subset X_i\cap X_j=\unsure{\bigcup|\coprod}_{k\leq i,j}X_k
\]
et le deuxième \unsure{mbr}, si $i\ne j$, \deleted{\unsure{et}} tout $k\leq i,j$ est ou bien $<i$, ou bien $<j$ (il ne peut être \unsure{égal} à la fois : $i$ et $j$) donc \deleted{\unsure{on}} $X_k\subset\partial X_i$ ou $X_k\subset\partial X_j$, donc
\[
i\ne j\Longrightarrow X_i\cap X_j\subset\partial X_i\cup\partial X_j\subset\complement(X_i^\circ\cap X_j^\circ)
\]
donc
\[
X_i^\circ\cap X_j^\circ\subset\complement X_i^\circ\cap X_j^\circ\qquad\text{d'où}\quad X_i^\circ\cap X_j^\circ=\varnothing
\]

Prouvons $3')$. \inserted{On a bien sûr $\supset$, il faut voir $\subset$.} Soit $x\in X_i$, par hyp. $\exists j\in I$ tel que $x\in X_j^\circ$, \unsure{tout} revient à voir que $X_j^\circ\subset X_i$.
\scanpage{51}
\mspage{42}
% archive sheet number: 51

i.e. la \unsure{sempiternelle} condition
\[
X_j^\circ\cap X_i\ne\varnothing\Longrightarrow j\leq i.
\]
Mais on a
\[
X_j^\circ\cap X_i\subset X_j\cap X_i=\bigcup_{k\leq i,j}X_k
\]
et si on n'avait pas $j\leq i$, on ne peut \unsure{avec} la somme du dernier mbr avoir $k=j$, donc cette somme est $\subset\bigcup_{k<j}X_k=\partial X_j\subset\complement X_j^\circ$ \inserted{$\subset\complement X_j^\circ\cap X_i$} donc $X_j^\circ\cap X_i$ est contenu dans son complémentaire, ce qui contredit l'hyp. que c'est vide.

La prop. est prouvée.

\uuline{Corollaire} $3$ Pour que $(X_i)_{i\in I}$ soit une figure \unsure{monoblaste} indexée par l'ens ordonné $I$ (cf Cor. $2$), il f et il s. qu'on ait les conditions suivantes
\[
\forall i,j\in I\ \text{avec }i\leq j,\ \text{on a }X_i\deleted{\unsure{\in}}\subset X_j\tag*{a)}
\]
\[
\forall i,j\in I,\ \text{on a }\deleted{X_i^\circ\ne\varnothing,\ \unsure{\Longleftrightarrow i\ne j}}\quad\inserted{i\ne j\Longrightarrow}\ X_i^\circ\cap X_j^\circ=\varnothing\tag*{\unsure{b)}}
\]
% slanted
\margin{\uline{NB} b équivaut à $2$ conditions $b')\ \forall i\in I,\ X_i^\circ\ne\varnothing$ $b'')\ \forall i,j,\ i\ne j\quad X_i^\circ\cap X_j^\circ=\varnothing$}

\unsure{c)} $\forall i\in I$, si on désigne par $\unsure{\Phi|P}_i\subset\unsure{\mathcal{P}|\mathrm{P}}(X)$ l'ens des parties de $X$ de la forme $X_j$, $j\leq i$, $\Phi_i$ est une figure \unsure{monoblaste}.

\deleted{on \unsure{trouve|prouvera}} \deleted{(\unsure{d'après} \illegible{2})}
\inserted{(Il suffit $=$ qu'elle soit ``quasi-admissible'', et \unsure{moins|même} encore, que \deleted{\unsure{$\forall$|tout}} $x\in X_i$, dans l'ens des $X_j$ ($j\leq i$) qui contiennent $x$, qu'il y ait un élément minimal (ce qui est \unsure{certain} p.ex. $I$ est fini !!))}

Tout revient à montrer que ces conditions \deleted{a) b)} impliquent \deleted{\unsure{pour le moins que}} \inserted{$3')\quad X_i=\bigcup_{j\leq i}X_j^\circ.$} (car on \deleted{\unsure{l'a}} \deleted{$\forall i,j\in I$} \unsure{a} \deleted{$2')$, et} $2')=b'')$.
\[
\deleted{X_i\subset X_j\Longleftrightarrow i\leq j}\qquad\deleted{(\Longleftarrow\text{ est connu par a})}
\]
% slanted
\margin{\uline{NB} les conditions b) \deleted{\unsure{c)}} impliquent déjà (c) via $\bigcup X_i^\circ=$}

\deleted{Ce qui impliquera : $2')$ et $3')$ donc (i).}

\deleted{Prouvons, $\forall i\in I$}
\[
\begin{aligned}
\deleted{X_i^\circ}&\deleted{=\text{cellule \unsure{ouverte|associée} \unsure{de|à} }X_i\text{ dans }\Phi_i.}\\
&\deleted{=\text{(on a donc) }X_i\setminus\bigcup_{\substack{j\leq i\\X_j\subsetneq X_i}}X_j}
\end{aligned}
\]
\[
\deleted{\unsure{\subset}}
\]
\scanpage{52}
\mspage{43}
% archive sheet number: 52

Fixons un $i\in I$. Pour tout $j\leq i$, posons
\[
\begin{aligned}
X_j^{\circ\prime}&=\text{cellule \unsure{ouverte} de }X_j\text{ dans }\unsure{\Phi|P}_i\\
&=X_j\setminus\bigcup_{\substack{j'\leq i\\\text{t.q. }X_{j'}\subsetneq X_j}}X_{j'}
\end{aligned}
\]
\margin{\uline{NB} dans le lemme, on ne suppose \uline{pas} $j'\leq j$}

\inserted{Je dis que} cette somme comporte plus de termes a priori que la somme analogue qui figure dans la déf. de
\[
X_j^\circ=X_j\setminus\bigcup_{j'<j}X_{j'},
\]
i.e. que
\[
j'<j\Longrightarrow X_{j'}\subsetneq X_j.
\]
En effet, on \deleted{\unsure{sait}} \inserted{a} \deleted{que $X_j$} $X_{j'}\subset\partial X_j$, donc \deleted{\unsure{on}} donc $X_j\setminus X_{j'}\supset X_j\setminus\partial X_j=X_j^\circ$, comme $X_j^\circ\ne\varnothing$ par hyp. on conclut \unsure{que} bien $X_j\setminus X_{j'}\ne\varnothing$ i.e. $X_{j'}\ne X_j$.

On a donc prouvé par là-même
\[
X_j^{\circ\prime}\subset X_j^\circ\qquad\text{``\unsure{bigre|bigger}''}
\]

\deleted{d'autre part par hyp. Mais l'ens des $X_j^{\circ\prime}$ \unsure{recouvre} $X_i$} \inserted{(\unsure{par la publ.} :)} \deleted{Or \unsure{a}} \deleted{Donc}
\[
\deleted{\bigcup_{j\leq i}X_j^{\circ\prime}\subset\bigcup_{j\leq i}X_j^\circ\subset X_i}
\]
% slanted
\margin{donc $\bigcup X_j^{\circ\prime}=X_i$ implique $\bigcup X_j^\circ=X_i$, qfd. \unsure{$\dagger$}}

\deleted{avec \uline{égalité} des deux extrêmes. Comme les $X_j^\circ$ sont $\ne\varnothing$ et deux à deux disjoints, cela implique}
\[
\deleted{X_j^{\circ\prime}=X_j^\circ:}
\]
\deleted{Les cellules ouvertes sont les $X_j^\circ$. En même} \deleted{En même temps,} \deleted{cela entraîne que} \deleted{temps, si $j,k\leq i$, je dis que}
\[
\deleted{\unsure{i}}\qquad\deleted{j\unsure{\nleq|\ne}k\Longrightarrow}
\]

\deleted{Soient $i,j\in I$ tels que $X_i\subset X_j$, prouvons $i\leq j$. Soit $x\in X_i^\circ$ ($X_i^\circ\ne\varnothing$ par $b')$}
\scanpage{53}
\mspage{44}
% archive sheet number: 53

\deleted{par c) $\exists j'\leq j$ avec $x\in X_{j'}^{\circ\prime}$, a fortiori $x\in X_{j'}^\circ$. Donc $X_i^\circ\cap X_{j'}^\circ\ne\varnothing$, donc par $b'')$, on a $i=j'$, d'où $i\leq j$, qfd.}

\uline{Exemple} Voici un exemple où on a $1'$ et $2'$, et en plus tous les $X_i^\circ\ne\varnothing$, sans avoir $3'$. \inserted{(indexation respectant la relation d'inclusion $i\leq j\Longleftrightarrow X_i\subset X_j$)}

Soit $\deleted{X=}A_0\supset A_1\supset\cdots\supset A_n\supset\cdots$ une suite \uline{strictement} \uline{décroissante} d'ens. de $X$, avec
\[
\deleted{A_\infty=}\inserted{N=}\bigcap A_n\ne\varnothing.
\]
\inserted{Soit $\deleted{\unsure{B}}A_\infty\subset X$ avec $A_\infty\cap A_0=N$, $A_\infty\not\subset A_0$ i.e. $N\ne A_\infty$.}
On prend $\deleted{\unsure{B\cap A_n}}\unsure{I|\Phi}\in\unsure{\mathcal{P}|\mathrm{P}}(X)$ défini comme ens. des $A_i$ ($i\in\mathbb{N}\cup\{\infty\}$). \deleted{\unsure{I est}} \inserted{$I\setminus\{\infty\}$} un ens \deleted{\unsure{tot.}} ordonné par inclusion, mais $\infty$ est un ``pt isolé''. On a
\[
\begin{aligned}
A_i^\circ&=A_i\setminus A_{i+1}\ne\varnothing&&\text{si }i\ne\infty\\
A_\infty^\circ&=A_\infty\ne\varnothing
\end{aligned}
\]
\deleted{\unsure{Les $A_i$}} Les $A_i^\circ$ ($i\in I$) sont mutuellement disjoints non vides, leur réunion est $A_0\inserted{\cup A_\infty}$ $(=|\unsure{\Phi|P}|)$. Mais Un $A_i$ ($i\ne\infty$) n'est pas réunion des $A_j^\circ$ avec $A_j\subset A_i$ (il manque $N$ !), il n'est pas saturé par la relation d'équivalence définie par les $A_j^\circ$ ($j\in I$).

\deleted{D'autre part,} \inserted{Un} $x\in N$ n'est pas contenu dans un plus petit élément de $\Phi$. On a
\[
\underline{A_\infty^\circ\cap A_i}\inserted{=N}\ne\varnothing,
\]
\margin{$\unsure{A_\infty^\circ\cap A_i}$}
sans pour autant avoir $A_\infty^\circ\subset A_i$ \deleted{(mais bien $A_\infty^\circ\subset A_i$}

Ici la condition \unsure{-bis $\beta$)} n'est pas satisfaite (saturation des $A_i$), alors que $\alpha$, $\gamma$) le sont.

Question : si on suppose
\[
\left\{\begin{aligned}
&\forall i\quad X_i^\circ\ne\varnothing\\
&\forall i\ne j\quad X_i^\circ\cap X_j^\circ=\varnothing\\
&\deleted{\bigcup X_i^\circ=|\Phi|}
\end{aligned}\right.
\qquad+\ \begin{gathered}\text{la saturation}\\\text{de chaque }X_i\end{gathered}
\]
\[
X_i=\bigcup_{j\in I,\ X_j^\circ\subset X_i}X_j
\]
\scanpage{54}
\mspage{45}
% archive sheet number: 54

a-t-on $\gamma)$ -- savoir que
\[
X_j^\circ\subset X_i\Longrightarrow j\leq i
\]
du tout ? Du moins si \deleted{l'indexation} on suppose
\[
X_j\unsure{\leq|\subseteq}X_i\Longrightarrow j\leq i\qquad\text{(``indexation respectant les inclusions'')}
\]
i.e.
\[
X_j^\circ\subset X_i\Longrightarrow\deleted{j\leq i}\ X_j\subset X_i\ ?
\]

Voici un \deleted{autre} contreexemple aussi à \unsure{cela|ceci}.

\deleted{Supposons qu'on ait \uline{deux} familles str. décroissantes, $(A_i)$ comme dessus et}
\[
\deleted{B_0\supset B_1\supset\cdots B_n\supset\cdots}
\]
\deleted{Supposons}
\[
\deleted{A_0\cap B_0=\bigcap A_i=\bigcap B_i,\quad\text{soit }N}
\]
\deleted{et prenons}
\[
\deleted{\unsure{\Phi|P}=\{A_i,B_j\}.}
\]
\deleted{On a $A_i$}

Rajoutons à la famille un dernier ensemble
\[
A=A_\infty-N
\]
i.e. prenons
\[
\Phi=\{(A_i),A_\infty,A\}
\]
On trouve
\[
A_i^\circ=A_i\setminus A_{i-1},\quad A_\infty^\circ=N,\ A^\circ=A
\]
tous ces ens sont $\ne\varnothing$ et mutuellement disjoints, réunion est $|\Phi|=A_0\cup A_\infty=A_0\unsure{\amalg|\sqcup}A$.
On a ici $A_\infty^\circ\subset A_i$ pour $i\ne\unsure{0|\infty}$, alors que $A_\infty\not\subset A_i$.
\scanpage{55}
\mspage{46}
% archive sheet number: 55

\uline{Homomorphismes \inserted{de \unsure{magasins} et} d'ateliers}

\deleted{Soient}

\uline{Ex 1} Soit $f:{\unsure{L|\mathscr{L}}}'\longrightarrow L$ une application ensembliste, on va en déduire une application canonique
\[
f^*:\mathrm{Figures}(L)\longrightarrow\mathrm{Figures}(L')
\]
Pour toute figure \inserted{\unsure{un|une} $\unsure{\Phi|P}$} dans $L$, on pose
\[
f^*(\Phi)=\{f^{-1}(A)\mid A\in\Phi,\ \underline{f^{-1}(A^\circ)\ne\varnothing}\}\tag*{(1)}
\]
\inserted{i.e. $A^\circ\cap\deleted{\unsure{f_0}}L_0\ne\varnothing$ où $L_0=f(L')$}

On peut considérer $f^*$ comme un composé de \deleted{deux} \inserted{trois} applications ; l'une
% slanted
\margin{\deleted{si $f$ est \unsure{surj.}} \unsure{voir}}
\[
A\longrightarrow\deleted{\unsure{de}}F_{L_0}\bullet A
\]
\margin{$F_{L_0}=\{L_0\}$, figure réduite à la seule strate $\{L_0\}$ si $L_0\ne\varnothing$ i.e. $L'\ne\varnothing$, sinon $F_{L_0}=\varnothing$}
va de
\[
\mathrm{Figures}(\unsure{L'|L})\longrightarrow\mathrm{Figures}(\unsure{L^\circ|L})_{L_0},\tag*{(1 bis)}
\]
le deuxième mbr désigne le sous-ens. du premier, formé des figures à \unsure{support} dans $L_0$. On a d'autre part
\[
\mathrm{Figures}_{L_0}(L)\xrightarrow{\sim}\mathrm{Figures}(L_0)\tag*{(1 ter)}
\]
% slanted
\margin{Le composé de (1 bis) et de (1 ter) est $i^*:\mathrm{Figures}(L)\longrightarrow\mathrm{Figures}(L_0)$ qui \unsure{formellement} \unsure{s'apparente} à un passage au quotient\ldots}
et enfin
\[
f_0^*:\mathrm{Figures}_{\deleted{\unsure{L}}}(L_0)\unsure{\relbar\joinrel\relbar|\longrightarrow}\mathrm{Figures}(L'),
\]
% slanted
\margin{$f_0$ étant l'appl. \unsure{can.} induite $\unsure{L|L'}\longrightarrow L_0$}
où cette fois, vu que $f$ est surjectif, on peut poser simplement
\[
f_0^*(\Phi_0)=\{f^{-1}(A)\mid A\in\Phi_0\},\tag*{(2)}
\]
la condition $f_0^{-1}(A^\circ)\ne\varnothing$ étant automatique
\scanpage{56}
\mspage{47}
% archive sheet number: 56

pour toute $A\in\unsure{\Phi|P}_0$.

L'application \unsure{(1)|(1 bis)} peut être vue comme un ``passage au quotient'', obtenu en posant dans $\mathrm{Fig}(\unsure{L|\mathscr{L}})$ formellement
\[
\text{``}X=\varnothing\text{''}
\]
pour tte multistrate $X$ telle que
\[
X^\circ\cap L_0=\varnothing\qquad\text{i.e.}\quad F_{L_0}\mathrel{|\circ|}X
\]
\deleted{Il y a \unsure{danger} \unsure{toutefois}, \unsure{une}}
(mais attention, on conserve les sous-strates $Y$ pour lesquelles $Y^\circ\cap L_0\ne\varnothing$).

\deleted{Il y} On devine que de façon générale, un tel passage au quotient aura un sens (quand on se donne, dans un magasin $M$, une partie qui soit $|\circ|$-fermée, ou un ``$|\circ|$''-support, dont les éléments seront \unsure{égalés} à $0$. Si on néglige $\deleted{\unsure{R}}X$, on ne néglige pas pour autant toutes ses strates, mais du moins tous les $Y$ tels que $Y\mathrel{\overset{\circ}{\ll}}X$ (\unsure{moralement}, tels que $Y^\circ\subset X^\circ$)
% slanted
\margin{NB}

\deleted{\unsure{L'} L'application}

L'application (1 ter) est un isomorphisme, inverse de
\diagram{\begin{tikzcd}[ampersand replacement=\&,column sep=small]
\mathrm{Figures}_{\deleted{\unsure{L_0}}}(L_0)\arrow[r,"\sim"] \& \mathrm{Figures}(L)_{L_0}\arrow[r,hook] \& \mathrm{Figures}(L)
\end{tikzcd}}
et le composé
\[
i_*:\mathrm{Figures}(L_0)\unsure{\relbar\joinrel\relbar|\longrightarrow}\mathrm{Figures}(L)
\]
est un ``plongement'', d'image $\mathrm{Figures}(L)_{L_0}$.
\scanpage{57}
\mspage{48}
% archive sheet number: 57

\unsure{Il} a l'air d'avoir des propriétés \unsure{similaires} \unsure{de|à} celles de $f_0^*$ (qui est elle aussi un plongement) avec une différence importante.

Dans les ateliers considérés, il y a un plus grand élément $F_{\unsure{L|\mathscr{L}}}$, $F_{L'}$, $F_{L_0}$. Dans les cas \inserted{(1),} \unsure{(2)} \inserted{(ou $\unsure{i_0^*}$)} et (1 ter), \deleted{\unsure{le}} (cas d'une application ``image inverse'' l'élément unité est transformé en \unsure{itou}, mais pas \inserted{en général} dans le cas d'une application ``image directe'' par une inclusion. Du pt de vue supports, on devine que la différence cruciale, c'est que dans le cas ${\unsure{\varphi|\phi}}^*$, le support \deleted{entier} total est transformé en \unsure{itou}, mais pas dans le cas $\varphi_*$. Et encore : $\varphi^*$ commute (sur les supports) au passage au complémentaire, et $\varphi_*$ pas.

Notons que dans le cas d'une inclusion $\varphi=i:L_0\hookrightarrow L$, \inserted{où} on aura à la fois un $i^*$, et un $i_*$, on aura
\[
i^*i_*=\mathrm{id}_{\unsure{F_0|\mathscr{F}_0}},\qquad\deleted{\unsure{i^*i_*}}i_*i^*(X)=i_*(1)\cdot X,\quad\text{plus gén}
\]
\[
i_*(X')\cdot Y=i_*(X'\cdot i^*(Y))
\]
(formule de projection)

À la fois les $\varphi^*$ et les $\varphi_*$ commutent aux $\mathrm{Sup}^{\leq}$ quelconques (y inclus $\varphi^*(\varnothing)=\varnothing$) et $\varphi_*$ aux $\mathrm{Inf}^{\leq}$ de familles non vides
